\section{Programación dinámica}

La \concept{programación dinámica} es un paradigma de diseño de algoritmos. Engañosamente, su nombre no guarda relación con cómo funciona.
%
\footnote{
El término "programación dinámica" fue acuñado por Richard Bellman, inventor del paradigma. En aquel entonces, el secretario de defensa estadounidense, Wilson, se oponía a la investigación y a la matemática, y la corporación RAND, donde trabajaba Bellman, estaba empleada por la fuerza aérea estadounidense. Bellman usó el término para describir sus actividades y protegerlas de Wilson: eligió la palabra "programación" porque su trabajo estaba relacionado con planeación y el adjetivo "dinámica" porque, según él, es imposible usarlo peyorativamente. (\cite{bellman_autobiography})
}%
. La idea del paradigma consiste en almacenar las soluciones a subproblemas que, de no almacernarse, tendrían que volverse a computar. El paradigma entonces acepta mayor complejidad espacial a cambio de menor complejidad temporal.

\subsubsection{Cuándo usar programación dinámica}

No tiene sentido intentar diseñar algoritmos de programación dinámica para todos los problemas. Los problemas que se resuelven eficientemente con programación dinámica \textit{siempre} presentan dos propiedades:
\begin{enumerate}
    \item \concept{Subestructura óptima}: la solución óptima del problema original se puede construir a partir de soluciones óptimas a sus subproblemas. Esto es un prerrequisito para que exista un algoritmo de programación dinámica.
    \item \concept{Subproblemas superpuestos}: varios subproblemas tienen subsubproblemas repetidos, de forma que un algoritmo recursivo acabaría resolviendo los mismos subsubproblemas varias veces.
    \begin{itemize}
        \item Para que eso ocurra, el espacio de posibles subproblemas debe ser pequeño, usualmente polinomial en el tamaño de la entrada. Si el espacio es muy grande, es probable que pocos o ningún subsubproblema se repita.
        \item Puede que exista un algoritmo de programación dinámica para el problema incluso si no presenta subproblemas superpuestos, pero entonces usar programación dinámica no es más eficiente que usar dividir y conquistar o alguna solución recursiva.
    \end{itemize}
\end{enumerate}

\subsubsection{Tipos de programación dinámica}

Hay en esencia dos formas de diseñar algoritmos de programación dinámica: 
\begin{itemize}
    \item \concept{\textit{Top-down}} (del problema más general a los más pequeños): Se resuelve el problema original de forma recursiva, pero cada vez que se resuelve un subproblema, se almacena su resultado para evitar tener que recomputarlo. A eso se le llama \concept{memoización}\footnote{El término "memoización" proviene del inglés \textit{memoization}, que significa convertir los resultados en un memo, abreviación del latín \textit{memorandum}: algo para recordar.}.
    \begin{itemize}
        \item Se crea una estructura de datos para almacenar los resultados de cada llamado recursivo. Inicialmente, se llena con valores que indiquen que el resultado en esa posición no se ha computado aún (e.g., \inline{null}, $\infty$, etc.). 
        
        Para que la información persista entre llamados recursivos, esa estructura tiene que estar definida afuera de la función recursiva; si no, cada llamado recursivo crearía su propia estructura, sin los resultados previos.
    \end{itemize}
    \item \concept{\textit{Bottom-up}} (de los problemas pequeños al general): Se resuelven primero los subproblemas más pequeños y se almacenan sus resultados en una tabla. A eso se le denomina \concept{tabulación}. Se usan esos resultados para resolver subproblemas más grandes, hasta llegar al problema original.
    \begin{itemize}
        \item \textbf{Iterativo}: Como los subproblemas se resuelven en orden de tamaño, se reemplaza el uso de recursión por un ciclo. 
        \item \textbf{Optimización de espacio}: Muchas soluciones bottom-up tienen la bondad de que, para tabular el siguiente término, no se necesitan todos los que se han tabulado hasta el momento, sino solo algunos. Se puede entonces sobreescribir el resto y reutilizar su espacio, reduciendo la complejidad espacial.
    \end{itemize}
\end{itemize}

Como ambos métodos aprovechan la propiedad de subproblemas superpuestos, sus complejidades asintóticas suelen ser iguales. Sin embargo, cuando \textit{todos} los subproblemas posibles se tienen que resolver al menos una vez, el algoritmo bottom-up suele ser más eficiente que el top-down por un factor constante, porque el algoritmo top-down tiene el \hyperref[sec:sobrecosto_recursion]{sobrecosto de la recursión}.

Ocasionalmente, no es necesario resolver todos los subproblemas del espacio, sino solo algunos; en ese caso, la solución bottom-up puede ser peor porque requiere resolver los subproblemas pequeños primero para completar la tabulación, mientras que la solución top-down solo memoiza los subproblemas que sí se tienen que resolver.

A continuación se presentan múltiples ejemplos de problemas que se resuelven eficientemente usando programación dinámica, para concretar estas ideas.

\setinputbase{diseno/programacion_dinamica}
\inputlist{
    fibonacci,
    corte_de_vara,
    % matrices_cadena,
    mochila
}
